\chapter{新玩具}

\begin{quote}
    我属我的良人，他也恋慕我。我的良人，来吧，你我可以往田间去。你我可以在村庄住宿！我们早晨起来往葡萄园去，看看葡萄发芽开花没有，石榴放蕊没有。我在那里要将我的爱情给你。

    \hfill《圣经·雅歌》7:11
\end{quote}

在本章，我会介绍专为本书创建的工具。为了看懂本章定义的大部分指令和环境，你需要阅读过第9章和第10章的内容……本章中，我们会涉及带有“危险”标记的提示的创建方法、各章开头的首字下沉、摘要、术语字典、带有当前章号的标签，以及并列展示\LaTeX 代码和效果的环境。

\section{一些小手工}

\subsection{参数和排版转换}

在涉及信息语言的文档中，需要突出显示参数、指令，以及函数。例如，我们需要这类效果：

\begin{codelist}[11.1]{
    为了编译文件\codereplace{文件}：
\begin{flushleft}
    \ttfamily latex \codereplace{文件}
\end{flushleft}
}\begin{verbatim}
为了编译文件\bwarg{文件}：
\begin{flushleft}
    \ttfamily latex \bwarg{文件}
\end{flushleft}
\end{verbatim}
\end{codelist}

指令\verb|\bwmarg|可以以\textsl{倾斜字体}来输出其参数，并且置于“⟨”和“⟩”之间。这种括号可以在数学模式下分别由指令\verb|\langle|、\verb|\rangle|生成。此外，你无疑注意到了，我们可以使用下标符号，就像这样：

\begin{codelist}[11.2]{
复制文件：
\begin{flushleft}\ttfamily
    cp \codereplace{文件$_1$} ...
    \codereplace{文件$_n$}
    \codereplace{文件$_{dst}$}
\end{flushleft} 
}\begin{verbatim}
复制文件：
\begin{flushleft}\ttfamily
    cp \bwarg[1]{文件} ...
    \bwarg[n]{文件}
    \bwarg[dst]{文件}
\end{flushleft}
\end{verbatim}
\end{codelist}

指令\verb|\bwarg|\jz{
    这个名称代表“黑白”参数（« black \& white » argument）……
}的定义如下：

\begin{dmd}
\begin{verbatim}
\newcommand{\marg}[2][]{% 
    {\normalfont%
        \textsl{$\langle$#2%
            % 如果选用了可选参数 
            \ifthenelse{\equal{#1}{}}{} 
            {$_\mathit{#1}$}% 显示下标
            $\rangle$}}}%
\end{verbatim}
\end{dmd}

指令\verb|\normalfont|的效果是回到文档的默认字体。这解释了为什么示例11.1中的“\codereplace{文件}”没有使用打字机字体显示。

在本文档的电子版（也就是需要通过屏幕阅读的版本）中，我们决定使用\emph{颜色}而不是字符“⟨”和“⟩”。这样一来，可以定义指令\verb|\colarg|：

\begin{dmd}
\begin{verbatim}
\newcommand{\colarg}[2][]{{% 
    \normalfont\color{blue!90}#2% 蓝色 
    \ifthenelse{\equal{#1}{}}{}{$_\mathit{#1}$}}}
\end{verbatim}
\end{dmd}

接下来，借助一个巧妙放置的布尔值\celan{\S 9.3.1}，我们可以定义一个通用的指令\verb|\marg|，以调用其中的一个版本（黑白版本或彩色版本）：

\begin{dmd}
\begin{verbatim}
\ifversionenligne
    \let\marg\colarg
\else
    \let\marg\bwarg
\fi
\end{verbatim}
\end{dmd}

这个结构调用了\TeX 的指令\verb|\let|，在9.2.3小节有明确的介绍。

\subsection{关于索引的生成}

本书中，凡是行文中提到指令、环境、包、文档类型等时，都会调用对应的特殊指令，来自动在索引中插入一条入口\yz{
    妈的怎么不早说。算了先不做索引了，全书翻译完之后再调整吧。
}。举例来说，这样一来的效果为：

\begin{codelist}[11.3]{
借助包\textsf{varioref}，可以
使用指令\dm{\backslash vref}……
}\begin{verbatim}
借助包\ltxpack{varioref}，可以
使用指令\ltxcom{vref}……
\end{verbatim}
\end{codelist}

指令\verb|\ltxpack|的定义如下。首先，以下内容定义的指令\verb|\ltx@pack|可以以非衬线字体展示包名：

\begin{dmd}
\begin{verbatim}
\newcommand{\ltx@pack}[1]{%
    \upshape\textsf{#1}}
\end{verbatim}
\end{dmd}

接下来，我们像这样定义：

\begin{dmd}
\begin{verbatim}
\newcommand{\ltxpack}[1]{%
    \ltx@pack{#1}% 
    \protect\index{扩展们！\protect\texttt{#1}}% 
    \protect\index{#1@\protect\textsf{#1 扩展}}}
\end{verbatim}
\end{dmd}

其中调用了前面定义的指令，并在索引中插入了两个入口，一个遵循“\codereplace{包名}扩展”的格式，另一个作为“扩展们”的\celan{\S 6.3}子入口。这里，指令\verb|\protect|的作用是避免指令\verb|\ltxpack|本身作为作为另一条指令的参数时带来的麻烦。遵循同样的思路，我们可以定义指令\verb|\ltxcom|。首先：

\begin{dmd}
\begin{verbatim}
\newcommand{\ltx@com}[1]{%
    \texttt{\symbol{92}#1}}
\end{verbatim}
\end{dmd}

以上指令可以以打字机字体生成指令名，并在前面加上字符\dm{\backslash}。\verb|\symbol|是\LaTeX 指令，此处用于插入所选字体下的第92个字符（恰好为反斜杠）。因此，我们最终有如下定义：

\begin{dmd}
\begin{verbatim}
\newcommand{\ltxcom}[1]{%
    \ltx@com{#1}% 
    \index{#1@\protect\texttt{\symbol{92}#1}}}
\end{verbatim}
\end{dmd}

该指令调用了上一条指令，在索引中插入一个入口。此处需要提炼出的思想是，定义可以自动在索引中插入入口的指令可能很有帮助。举例来说，我们也许可以定义这样一条指令：

\begin{dmd}
\begin{verbatim}
\newcommand{\jargonanglais}[1]{% 
    \emph{#1}%
    \index{#1}}
\end{verbatim}
\end{dmd}

该指令可以在设置英文术语格式的同时将其插入索引——比如某个特殊的索引。同样，如果在文档中经常出现，我们就可以定义一条指令，将其插到索引中。例如，在本书中，我们有这样的定义：

\begin{dmd}
\begin{verbatim}
\newcommand{\postscript}{% 
    PostScript% 
    \protect\index{PostScript}}
\end{verbatim}
\end{dmd}

\subsection{跳转提示}

本书的纸质版中四处分散着跳转阅读的提示，例如这个提醒你去看\celan{\S D.3.5}术语字典（glossaire）的提示——幸亏我们这里在说跳转阅读，否则这个术语字典就这里的内容毫无关联了\jz{
    如果你没完全跟上的话，我是想说，这里其实没有术语字典的什么事……
}。实现跳转提示符号的指令被赐名为\verb|\voir|，接收两个参数：

\begin{dmd}
\backslash voir\{\codereplace{目标标签}\}\{\codereplace{跳转对象文本}\}
\end{dmd}

例如，前面的跳转是这样生成的：

\begin{dmd}
\verb|\voir{chap-glossaire}{glossaire}|
\end{dmd}

这个指令的设计过程中，全部的“难点”在于怎样让三角形根据页面奇偶性来改变朝向。借助包\textsf{chngpage}，这一难点得以功课。请参阅9.3.3小节。

余下的工作就是展示三角形。为了达成这一目的，我们定义了两个指令，分别生成侧栏的跳转提示和正文中的标记。由此，\verb|\voir|的格式如下：

\begin{dmd}
\begin{verbatim}
\newcommand{\voir}[3][\S]{% 
    \checkoddpage% 
    \ifcpoddpage
        \v@irpageimpaire{#1}{#2}{#3}}{% 奇数页的跳转提示
    \else
        \v@irpagepaire{#1}{#2}{#3}}} % 偶数页的跳转提示
    \fi
\end{verbatim}
\end{dmd}%TODO 此处括号似乎不配对

我们注意到，除了两个必需的参数外，该指令还可以接收一个可选参数，它默认定义为段落标记（\S）。指令\verb|\v@irpageimpaire|和\verb|\v@irpagepaire|是对称的，作用如下：

\begin{enumerate}
    \item 将“朝向正确”的三角形放置在文本中，作为跳转提示对象；
    \item 在侧栏生成带有跳转目标的注解。
\end{enumerate}

三角形可以从包\textsf{amssymb}中包含的符号中获取：

\begin{codelist}[11.4]{
    天啊，那些“美观”的三角形：
$\blacktriangleleft$和
$\blacktriangleright$！
}\begin{verbatim}
天啊，那些“蟀爆了”的三角形：
$\blacktriangleleft$和
$\blacktriangleright$！
\end{verbatim}
\end{codelist}

如下所示，是针对在偶数页实现跳转提示的指令的最终版本：

\begin{dmd}
\begin{verbatim}
\newcommand{\v@irpageimpaire}[3]{% 
    {\tiny$\blacktriangleright$}#3% 文本中的跳转提示
    \marginpar{%
        \parbox[t]{.9\marginparwidth}{% 
            {\footnotesize\sffamily%
                \hfill#1~\ref{#2}}~{\small$\blacktriangleleft$}\\
            {\footnotesize\sffamily%
                \mbox{}\hfill p.~\pageref{#2}\hfill\mbox{}}}}
\end{verbatim}
\end{dmd}% TODO 括号配对

注意到，在侧栏的跳转内容包含于一个\verb|\parbox|\celan{\S 4.4.3}，其中分两行展示章节编号和页码\yz{
    中文版暂时没做页码。%TODO 页码
}。对于奇数页，需要做与偶数页相同的事情，只不过需要考虑到它是奇数页\dm{:-)}

在“在线”版本中，跳转提示可以像超文本链接（lien hypertexte）那样呈现。这可以简单地通过指令\verb|\hyperref|（对应包同名）来实现。因此，我们可以编写出这类指令：

\begin{dmd}
\begin{verbatim}
\newcommand{\voir}[3][\S]{% 
    \hyperref[#2]{#3}}
\end{verbatim}
\end{dmd}

其中的可选参数没有过多的用处，只是为了与“纸质”版本中的指令\verb+\voir+兼容。

\subsection{更改侧栏}

在本文中，我多次临时改变了侧栏的样子，尤其是针对用于并列展示\LaTeX 代码和其效果，或需要展示题记等内容%TODO 是题记吗？
的情况。为了实现这种效果，维护\LaTeX 法文问答社区的玛丽-保罗·克卢特（Marie-Paul Kluth）\jz{
    这位似乎有个命中注定的姓［lom prédestilé；译注：将nom prédestiné中的n替换为了l。此处是在调侃此人的姓（Kluth）与发明\TeX 的克努特（Knuth）仅有一个字母的区别］……
}提议使用一种集成环境，如下所示：

\begin{dmd}
\begin{verbatim}
\newenvironment{changemargin}[2]% 
    {\begin{list}{}{%
        \setlength{\listparindent}{\parindent}% 
        \setlength{\itemindent}{\parindent}% 
        \setlength{\leftmargin}{#1}% 
        \setlength{\rightmargin}{#2}%
    }\item }% 
    {\end{list}}
\end{verbatim}
\end{dmd}

其中的思路是定义一个列表，而我们修改其中的侧栏。该环境接收的两个参数分别代表左侧和右侧的边距尺寸。这样的思路可以让环境产生很有趣的结果：边距会根据页面的奇偶性而实用不同的尺寸。这样的环境可以定义如下：

\begin{dmd}
\begin{verbatim}
\newenvironment{agrandirmarges}[2]{% 
    \begin{list}{}{%
        \setlength{\topsep}{0pt}% 
        \setlength{\listparindent}{\parindent}% 
        \setlength{\itemindent}{\parindent}% 
        \setlength{\parsep}{0pt plus 1pt}% 
        \checkoddpage%
        \ifcpoddpage
            \setlength{\leftmargin}{-#1}
            \setlength{\rightmargin}{-#2}
        \else
            \setlength{\leftmargin}{-#2}
            \setlength{\rightmargin}{-#1}
        \fi}\item }% 
    {\end{list}}
\end{verbatim}
\end{dmd}

\newenvironment{agrandirmarges}[2]{% 
    \begin{list}{}{%
        \setlength{\topsep}{0pt}% 
        \setlength{\listparindent}{\parindent}% 
        \setlength{\itemindent}{\parindent}% 
        \setlength{\parsep}{0pt plus 1pt}% 
        \checkoddpage%
        \ifcpoddpage
            \setlength{\leftmargin}{-#1}
            \setlength{\rightmargin}{-#2}
        \else
            \setlength{\leftmargin}{-#2}
            \setlength{\rightmargin}{-#1}
        \fi}\item }% 
    {\end{list}}

注意，这里我们使用了指令\verb|\isodd|来测试页面的奇偶性。图\ref{fig:11.1}展示了使用这种环境的示例，其代码如下：

\begin{figure}[bt]
    \centering
    \newcommand{\graphtex}{%
      \includegraphics[height=2cm]{img/knuth-tex}}
    \begin{agrandirmarges}{1cm}{2cm}
      \begin{center}
        \graphtex\graphtex\graphtex\graphtex\graphtex%
        \graphtex\graphtex\graphtex\graphtex\graphtex%
        \graphtex\graphtex\graphtex\graphtex\graphtex%
        \graphtex\graphtex\graphtex\graphtex%
      \end{center}
    \caption{一幅内容完全不重要的图。它的出现只是为了说明我们可以在需要额外空间的时候临时修改左边距和右边距……}
    \label{fig:11.1}
    \end{agrandirmarges}
\end{figure}

\begin{dmd}
\begin{verbatim}
    \begin{figure}[tb]
        \begin{agrandirmarges}{1cm}{2cm}
            % 这里1cm针对“订口”侧
            %     2cm针对“切口”侧
        ...
        \caption{一幅内容完全不重要的……}
        \end{agrandirmarges}
      \end{figure}
\end{verbatim}
\end{dmd}

\section{各种提示}

提示中可以见到的符号来源于“剪贴画”（……）。在图\ref{fig:11.2}中以3 cm的宽度展示了这些符号。在文档对应位置插入这些符号的“提示”由鄙人定义的环境生成。环境基于\TeX 层面的一项功能，而这项功能正是我勤勉地阅读\TeX Book是发现的，那就是指令\verb|\parshape|。该指令可以为段落指定任意形状：

\begin{codelist}[11.5]{
\parshape=5
1.5cm 2.5cm
1.5cm 2.5cm
1cm 2cm
1cm 2cm
0pt \textwidth 
a a a a a a a a a a a a a a a a a a a 
a a a a a a a a a a a a a a a a a a a 
a a a a a a a a a a a a a a a a a a a 
}\begin{verbatim}
\parshape=5
1.5cm 2.5cm
1.5cm 2.5cm
1cm 2cm
1cm 2cm
0pt \textwidth 
a a a a a a a a a a a a a a a a a a a 
a a a a a a a a a a a a a a a a a a a 
a a a a a a a a a a a a a a a a a a a 
\end{verbatim}
\end{codelist}

\begin{figure}[htb]
    \begin{center}
      \includegraphics[width=3cm]{img/fmb-important}\quad%
      \includegraphics[width=3cm]{img/fmb-note}\quad%
      \includegraphics[width=3cm]{img/fmb-question}
    \end{center}
    \caption{手册中使用的符号}
    \label{fig:11.2}
\end{figure}

符号“=”后面的数字用于指明我们想要变形的行数。接下来，成对出现的数字用于指明缩进和该行变形后的长度。对于上述示例，有：

\begin{itemize}
    \item 最先出现的两行缩进1.5 cm，并且每行长2.5 cm；
    \item 接下来的两行缩进1 cm，并且每行长2 cm；
    \item 第五行（也是最后一次指明内容）可以决定接下来所有行的风格，即缩进0 cm且行长等于事先定义的长度\verb|\textwidth|。
\end{itemize}

因此，想要在段落中插入提示，我们可以借助该指令将前两行错位：

\newlength{\larnota}
\newlength{\largligne}

\begin{codelist}[11.6]{
\setlength{\larnota}{.9cm}
\setlength{\largligne}{%
    \textwidth-\larnota}
\parshape=3
\larnota\largligne
\larnota\largligne
0pt\textwidth
\noindent 注意，这个段落唯一的
作用，就是用来证明我们确实可以
将段落中的两行内容错后一些，并
让随后的内容如同什么都没有发生
一样继续……
}\begin{verbatim}
\setlength{\larnota}{.9cm}
\setlength{\largligne}{%
    \textwidth-\larnota}
\parshape=3
\larnota\largligne
\larnota\largligne
0pt\textwidth
\noindent 注意，这个段落唯一的
作用，就是用来证明我们确实可以
将段落中的两行内容错后一些，并
让随后的内容如同什么都没有发生
一样继续……
\end{verbatim}
\end{codelist}

好了，接下来要做的事情就是在指令\verb|\parshape|留下的“洞口”中放上图片，让我们来简单地尝试一下：

\begin{codelist}[11.7]{
\setlength{\larnota}{.9cm}
\setlength{\largligne}{%
    \textwidth-\larnota}
\parshape=3
\larnota\largligne\larnota\largligne
0pt\textwidth\noindent%
\includegraphics[width=\larnota]{%
img/fmb-important}
注意，这个段落唯一的作用，就是
用来证明我们确实可以将段落中的
两行内容错后一些，并让随后的内
容如同什么都没有发生一样继续……
}\begin{verbatim}
\setlength{\larnota}{.9cm}
\setlength{\largligne}{%
    \textwidth-\larnota}
\parshape=3
\larnota\largligne\larnota\largligne
0pt\textwidth\noindent%
\includegraphics[width=\larnota]{%
    \ficnota}
注意，这个段落唯一的作用，就是
用来证明我们确实可以将段落中的
两行内容错后一些，并让随后的内
容如同什么都没有发生一样继续……
\end{verbatim}
\end{codelist}

显然，这里的图像像其他字符一样放置在行上。我们将其放入宽度为零的字盒中，这样可以让字盒的内容右对齐：

\begin{codelist}[11.7]{
\setlength{\larnota}{.9cm}
\setlength{\largligne}{%
    \textwidth-\larnota}%
\parshape=3
\larnota\largligne\larnota\largligne%
0pt\textwidth\noindent%
\makebox[0pt][r]{%
    \includegraphics[width=\larnota]{%
    img/fmb-important}}%
注意，这个美观的段落唯一的作用
，就是用来证明我们确实可以将段
落中的两行内容错后一些 [...]
}\begin{verbatim}
\setlength{\larnota}{.9cm}
\setlength{\largligne}{%
    \textwidth-\larnota}%
\parshape=3
\larnota\largligne\larnota\largligne%
0pt\textwidth\noindent%
\makebox[0pt][r]{%
    \includegraphics[width=\larnota]{%
        \ficnota}}%
注意，这个美观的段落唯一的作用
，就是用来证明我们确实可以将段
落中的两行内容错后一些 [...]
\end{verbatim}
\end{codelist}

接下来要做的是为这个图标施加竖直方向的位移（实际上，这个图标是正方形的，因此我们可以使用尺寸\verb|\indnota|）：

\begin{codelist}[11.9]{
\setlength{\larnota}{.9cm}
\setlength{\largligne}{%
    \textwidth-\larnota}
\parshape=3
\larnota\largligne\larnota\largligne
0pt\textwidth\noindent%
\raisebox{-\larnota}{%
    \makebox[0pt][r]{%
        \includegraphics[width=\larnota]{%
        img/fmb-important}}}%
注意，这个美观的段落唯一的作用
，就是用来证明我们确实可以将段
落中的两行内容错后一些 [...]
}\begin{verbatim}
\setlength{\larnota}{.9cm}
\setlength{\largligne}{%
    \textwidth-\larnota}
\parshape=3
\larnota\largligne\larnota\largligne
0pt\textwidth\noindent%
\raisebox{-\larnota}{%
    \makebox[0pt][r]{%
        \includegraphics[width=\larnota]{%
            \ficnota}}}%
注意，这个美观的段落唯一的作用
，就是用来证明我们确实可以将段
落中的两行内容错后一些 [...]
\end{verbatim}
\end{codelist}

得，还是不行。需要让\LaTeX 相信我们所移动的字盒的尺寸为零：

\begin{codelist}[11.10]{
\setlength{\larnota}{.9cm}
\setlength{\largligne}{%
    \textwidth-\larnota}
\parshape=3
\larnota\largligne\larnota\largligne
0pt\textwidth\noindent%
\raisebox{-\larnota}[0pt][0pt]{%
    \makebox[0pt][r]{%
        \includegraphics[width=\larnota]{%
        img/fmb-important}}}%
注意，这个美观的段落唯一的作用
，就是用来证明我们确实可以将段
落中的两行内容错后一些 [...]
}\begin{verbatim}
\setlength{\larnota}{.9cm}
\setlength{\largligne}{%
    \textwidth-\larnota}
\parshape=3
\larnota\largligne\larnota\largligne
0pt\textwidth\noindent%
\raisebox{-\larnota}[0pt][0pt]{%
    \makebox[0pt][r]{%
        \includegraphics[width=\larnota]{%
            \ficnota}}}%
注意，这个美观的段落唯一的作用
，就是用来证明我们确实可以将段
落中的两行内容错后一些 [...]
\end{verbatim}
\end{codelist}

我们就快要达到目的了。接下来需要做两处调整。

\begin{itemize}
    \item 字盒太低了些，这是因为参考线是文本行的底端。因此，我们可以在位移时抬升\dm{1ex}（字符的高度）。
    \item 在图标和文档之间添加一些空间是个不错的想法。为此，我们定义一个长度\verb|\padnota|。
\end{itemize}

\newlength{\padnota}
\newlength{\indnota}

\begin{codelist}[11.11]{
\setlength{\padnota}{5pt}
\setlength{\larnota}{.9cm}
\setlength{\indnota}{\larnota+\padnota}
\setlength{\largligne}{%
    \textwidth-\indnota}
\parshape=3
\indnota\largligne\indnota\largligne
0pt\textwidth\noindent%
\raisebox{-\larnota+2.2ex}[0pt][0pt]{%
    \makebox[0pt][r]{%
    \includegraphics[width=\larnota]{%
    img/fmb-important}%
    \hspace*{\padnota}}}%
注意，这个美观的段落唯一的作用
，就是用来证明我们确实可以将段
落中的两行内容错后一些 [...]
}\begin{verbatim}
\setlength{\padnota}{5pt}
\setlength{\larnota}{.9cm}
\setlength{\indnota}{\larnota+\padnota}
\setlength{\largligne}{%
    \textwidth-\indnota}
\parshape=3
\indnota\largligne\indnota\largligne
0pt\textwidth\noindent%
\raisebox{-\larnota+2.2ex}[0pt][0pt]{%
    \makebox[0pt][r]{%
    \includegraphics[width=\larnota]{%
        \ficnota}%
    \hspace*{\padnota}}}%
注意，这个美观的段落唯一的作用
，就是用来证明我们确实可以将段
落中的两行内容错后一些 [...]
\end{verbatim}
\end{codelist}

当当当当！接下来将这些代码修改成可以创建新环境的样子就可以了。这里使用的技巧基于9.5节介绍的环境\dm{list}。环境的完整代码如下：

\begin{dmd}
\begin{verbatim}
% 为参数传入文件名
\newenvironment{pictonote}[1]{% 
    \begin{list}{}{%
        \setlength{\labelsep}{0pt}%
        \setlength{\rightmargin}{15pt}} 
    \item%
        \setlength{\indentationnota}{% 
            \@totalleftmargin+\largeurnota+\paddingnota}%
        \setlength{\largeurlignenota}{% 
            \linewidth-\largeurnota-\paddingnota}%
        \parshape=3% 
        \indentationnota\largeurlignenota% 
        \indentationnota\largeurlignenota% 
        \@totalleftmargin\linewidth% 
        \raisebox{-\largeurnota+2.2ex}[0pt][0pt]{%
            \makebox[0pt][r]{% 
                \includegraphics[width=\largeurnota]{#1}%
                 \hspace{\paddingnota}}}%
        \ignorespaces}{% 
    \end{list}}
\end{verbatim}
\end{dmd}

\begin{exclamation}
注意，使用尺寸\verb|\@totalleftmargin|可以获取列表左侧边距的宽度，并将可能出现的嵌套纳入考虑。实际上，列表中的尺寸\verb|\leftmargin|代表了\emph{相对于}包含该列表的环境的做边距。
\end{exclamation}

接下来，我们可以使用此环境，并为其传入特殊的文件作为参数。更好的方案是，我们可以去定义新的环境，例如：

\begin{dmd}
\verb+\newenvironment{nota}{%+
\verb+    \begin{pictonote}{+\codereplace{你的文件}\verb+}}{\end{pictnote}}+
\end{dmd}

\begin{qquestion}
作为本节的结尾，需要知道，这个环境有些问题：如果提示中包含新的段落，那么会多出一块为图标预留的空白。%TODO

没错，没错，就像这里展示的一样。你也看到了，段首为图标留出了一块空白，但这里我们并不需要，对吧？为了避免这种麻烦，需要借助下面的咒语，在\TeX 的层面重新初始化缩进：

\begin{dmd}
\verb|\par\parshape=1\@totalleftmargin\linewidth|
\end{dmd}

这段内容可以作为一个指令的对象，在各个新段落的段首手动插入，就像我刚刚做的那样\dm{:-)}
\end{qquestion}

\section{引文}

\subsection{题记}

